\documentclass[hidelinks]{article}
\usepackage[a4paper, total={7in, 10in}]{geometry}
\usepackage[dvipsnames]{xcolor}
\usepackage{amsmath, amssymb, amsthm}
\usepackage{tikz}
\usepackage{tkz-euclide}
\usepackage[unicode]{hyperref}
\usepackage[all]{hypcap}
\usepackage{fancyhdr}
\usepackage{listings}
\usepackage{subcaption}
\usepackage{enumitem}
\usepackage{graphicx}

\lstset{basicstyle=\ttfamily\small,columns=fullflexible,frame=single,breaklines=true,showstringspaces=false}
\usetikzlibrary{angles,calc,decorations.pathreplacing}

\definecolor{carminered}{rgb}{1.0, 0.0, 0.22}
\definecolor{capri}{rgb}{0.0, 0.75, 1.0}
\definecolor{brightlavender}{rgb}{0.75, 0.58, 0.89}

\title{\textbf{Poisson's Equation Test}}
\author{Diego Juarez}
\date{March 25th, 2026}

\pagestyle{fancy}
\fancyhf{}
\fancyhead[L]{Poisson's Equation Test}
\fancyhead[R]{Jackson-style Practice}
\fancyfoot[C]{\thepage}

\begin{document}
\hypersetup{bookmarksnumbered=true}
\pagecolor{black}
\color{white}
\maketitle

\begin{Large}
\tableofcontents
\end{Large}
\pagebreak

\section{Source Notebook}
This problem set follows the notebook
\url{../notebooks/Poisson's equation.ipynb}
and focuses on electrostatic source problems, boundary conditions, and symmetry reductions.

\section{Core Theory}
\subsection{Problem 1: From Maxwell to Poisson}
Starting from Gauss's law in differential form,
\[
\nabla\cdot \mathbf{E}=\frac{\rho}{\varepsilon_0},
\]
and using $\mathbf{E}=-\nabla \Phi$, derive Poisson's equation. State clearly when it reduces to Laplace's equation.

\subsection{Problem 2: Linearity}
Suppose $\Phi_1$ and $\Phi_2$ satisfy
\[
\nabla^2\Phi_1=-\frac{\rho_1}{\varepsilon_0}, \qquad
\nabla^2\Phi_2=-\frac{\rho_2}{\varepsilon_0}.
\]
Show that any linear combination $a\Phi_1+b\Phi_2$ satisfies a corresponding Poisson equation. Explain the physical meaning of this result.

\subsection{Problem 3: Units}
Check the SI units of each term in
\[
\nabla^2\Phi=-\frac{\rho}{\varepsilon_0}.
\]
Confirm that the equation is dimensionally consistent.

\section{Canonical Examples}
\subsection{Problem 4: Uniform Charge Density in a Slab}
An infinite slab occupies $-a<z<a$ and has constant volume charge density $\rho_0$.
\begin{enumerate}[label=\alph*)]
\item Derive the one-dimensional Poisson equation for $\Phi(z)$.
\item Solve for the electric field inside and outside the slab using symmetry.
\item Obtain the potential, taking $\Phi(0)=0$.
\item Check continuity conditions at $z=\pm a$.
\end{enumerate}

\subsection{Problem 5: Point Charge in Three Dimensions}
Show that for a point charge $q$ at the origin,
\[
\rho(\mathbf{r})=q\,\delta^{(3)}(\mathbf{r}),
\]
Poisson's equation becomes
\[
\nabla^2\Phi(\mathbf{r})=-\frac{q}{\varepsilon_0}\delta^{(3)}(\mathbf{r}).
\]
Verify that
\[
\Phi(\mathbf{r})=\frac{1}{4\pi\varepsilon_0}\frac{q}{r}
\]
solves the equation away from the origin, and explain the role of the singularity at $r=0$.

\subsection{Problem 6: Spherical Symmetry}
Assume a spherically symmetric potential $\Phi(r)$.
\begin{enumerate}[label=\alph*)]
\item Write Poisson's equation in spherical symmetry.
\item Solve it for a general radial charge density $\rho(r)$ in integral form.
\item State what simplification occurs in regions where $\rho(r)=0$.
\end{enumerate}

\subsection{Problem 7: Uniformly Charged Solid Sphere}
A sphere of radius $R$ has uniform charge density $\rho_0$.
\begin{enumerate}[label=\alph*)]
\item Find the electric field inside and outside the sphere.
\item Derive the potential everywhere, imposing $\Phi(\infty)=0$.
\item Check continuity of $\Phi$ and $E_r$ at $r=R$.
\end{enumerate}

\subsection{Problem 8: Piecewise One-Dimensional Source}
Consider
\[
\rho(x)=
\begin{cases}
\rho_0, & |x|<a,\\
0, & |x|>a.
\end{cases}
\]
Solve the one-dimensional Poisson equation for $\Phi(x)$ piecewise. Determine the integration constants from continuity and symmetry assumptions.

\section{Boundary Value Viewpoint}
\subsection{Problem 9: Uniqueness}
State and prove a uniqueness theorem for Poisson's equation with Dirichlet boundary conditions. Your proof may use the energy method.

\subsection{Problem 10: Boundary Conditions}
At an interface between two regions with different charge densities, state the continuity properties required of $\Phi$ and of the normal derivative of $\Phi$. Explain how a surface charge modifies the normal derivative condition.

\subsection{Problem 11: Region with Zero Charge}
Suppose $\rho=0$ in a region $V$, but $\rho\neq 0$ elsewhere. Explain why the potential inside $V$ still depends on charges outside the region. Give a concrete electrostatic example.

\section{Green-Function Preview}
\subsection{Problem 12: Integral Representation}
Assuming free space and suitable decay at infinity, show formally that the solution of Poisson's equation can be written as
\[
\Phi(\mathbf{r})=\frac{1}{4\pi\varepsilon_0}\int \frac{\rho(\mathbf{r}')}{|\mathbf{r}-\mathbf{r}'|}\,d^3r'.
\]
Explain how this connects Poisson's equation with Coulomb's law.

\section{Computational Interpretation}
\subsection{Problem 13: Curvature and Charge}
In one dimension,
\[
\frac{d^2\Phi}{dx^2}=-\frac{\rho(x)}{\varepsilon_0}.
\]
Explain the geometric relation between the sign of $\rho(x)$ and the curvature of $\Phi(x)$. What does a localized positive charge do to the graph of the potential?

\subsection{Problem 14: Finite-Difference Form}
On a uniform one-dimensional grid with spacing $h$, derive the finite-difference approximation
\[
\frac{\Phi_{i+1}-2\Phi_i+\Phi_{i-1}}{h^2}=-\frac{\rho_i}{\varepsilon_0}.
\]
State the order of the truncation error and explain how boundary conditions enter the discrete system.

\section{Challenge Problems}
\subsection{Problem 15: Mean-Zero Compatibility on a Finite Interval}
Consider
\[
\Phi''(x)=-f(x), \qquad 0<x<L,
\]
with Neumann conditions $\Phi'(0)=\Phi'(L)=0$.
\begin{enumerate}[label=\alph*)]
\item Show that a solution can exist only if
\[
\int_0^L f(x)\,dx=0.
\]
\item Interpret this condition physically.
\end{enumerate}

\subsection{Problem 16: Delta-Function Source in 1D}
Solve
\[
\Phi''(x)=-\frac{q}{\varepsilon_0}\delta(x-a)
\]
on the interval $0<x<L$ with boundary conditions $\Phi(0)=\Phi(L)=0$.
State the continuity condition on $\Phi$ and the jump condition on $\Phi'$ at $x=a$.

\end{document}
